% MT053 / explanation — canonical editable content.
先区分这里的两个“后验”。\(P(y\mid x)\)是在知道输入后，\textbf{类别}有多可能；\(p(w\mid D)\)是在知道数据集\(D\)后，\textbf{参数}有多可能。LR直接描述前者，给权重加先验并做MAP时才涉及后者。



\begin{itemize}

\item \textbf{A错误。} 闭式解是能通过有限公式直接算出最优参数，例如某些线性最小二乘问题。一般LR的梯度中含依赖\(w\)的sigmoid，把梯度设为0后不能直接解出一个通用矩阵公式，通常要迭代数值优化。

\item \textbf{B正确。} 高斯先验偏好不太大的权重；取负log后变成L2惩罚，可以控制过拟合。这个先验加在参数\(w\)上，与生成式模型的类别先验\(P(y)\)不是同一项。

\item \textbf{C正确。} LR直接建模\(P(y\mid x)\)；生成式学习\(p(x\mid y)\)与\(P(y)\)，再导出后验。

\item \textbf{D正确。} 正则化强度是超参数，可用训练数据内部的CV选择，使限制强度适合新样本。

\item \textbf{E错误。} 对\(\alpha\|w\|^2\)求导得到\(2\alpha w\)，梯度更新因此多了一项把权重拉向0的作用，最优系数会改变。

\end{itemize}


\textbf{出处：} 2023B期中 Q2，原卷第2页（5分）


\textbf{答案依据：} 2023B随卷解答 Q2，随卷答案BCD；其余选项为跨讲说明。
